\section{Related Work}
Four areas inform the controller developed here: UAV-assisted vehicular
edge computing, queue-aware stochastic control, multi-agent reinforcement
learning, and robust policy training. We organize the discussion around the
capabilities each area contributes to the task-service and control problem.

\subsection{UAV-Assisted MEC and End-to-End Task Service}
UAV-mounted edge servers can shorten access distance and follow spatially
varying demand, but their trajectory, radio service, computation, and energy
budgets are coupled. Joint offloading, CPU control, wireless power delivery,
and UAV motion have therefore been studied as one optimization problem
\cite{liu2020uav,xiong2021energy}. Multi-UAV work further considers
cooperative offloading and resource allocation under shared airspace and
capacity constraints \cite{guo2024multi}, while recent vehicular MEC models
use multi-agent policies for task partitioning and resource allocation
\cite{liang2025rav}. These studies motivate joint control, but a task is often
represented by a single delay expression or one-slot allocation.

Our model records the stages a task actually visits: ingress, radio upload,
receive staging, edge computation, UAV--cloud backhaul, cloud computation,
and branch synchronization. Retaining these stages makes it possible to
account for waiting beyond the radio link. Tasks that miss their deadlines
remain in the service model until completion or the end of the trace.

\subsection{Queue-Aware Vehicular Offloading and Lyapunov Control}
Lyapunov optimization converts a long-horizon stochastic control problem into
queue-weighted slot decisions when arrivals, service, and energy increments
satisfy the required boundedness and feasibility conditions
\cite{neely2010stochastic}. It has been applied to MEC offloading with energy
harvesting \cite{mao2017dynamic}. Q-TOP combines physical queues with
three-dimensional multi-UAV path planning for vehicular edge computing and
uses a genetic optimizer for the resulting per-slot problem
\cite{sun2026qtop}.

The proposed formulation retains this queue-sensitive structure but adds
explicit transport, staging, backhaul, and cloud queues, task identities, and
a virtual net-energy deficit. Moving work from one stage to another is not
treated as completion: the drift term accounts for congestion at both the
source and receiving queues.

\subsection{Multi-Agent and Hybrid-Action Reinforcement Learning}
MADDPG uses centralized critics during training and decentralized actors at
execution, allowing agents with different observations to optimize a common
objective \cite{lowe2017multi}. MAPPO provides an on-policy alternative with
clipped policy updates \cite{schulman2017proximal}. Attention-based MAPPO has
also been used to coordinate bounded UAV motion, task service, and energy
transfer \cite{talukdar2026mappo}, while generative-policy extensions have
been studied for UAV task assignment \cite{tang2025dnn}.

The action space considered here is mixed. Vehicles select a categorical
hold or destination decision and a continuous partition fraction; RSUs and
UAVs allocate continuous processor shares; and UAVs additionally choose CPU
frequency and three-dimensional motion. The controller therefore combines
masked straight-through categorical decisions, bounded continuous heads,
and a deterministic feasibility layer. Neither MADDPG nor PPO alone supplies
battery causality, queue stability, or collision avoidance.

\subsection{Meta-Learning and Robustness Across Operating Domains}
MAML and Reptile learn an initialization intended to adapt with a small
number of gradient updates \cite{finn2017model,nichol2018reptile}. Yan
\emph{et al.} combine Reptile and MADDPG for computation allocation in a
UAV/RSU/cloud MEC architecture \cite{yan2026rmaddpg}. Domain randomization
addresses a different objective: one policy is optimized over a distribution
of environments and then deployed without online parameter updates. This
principle has been used to improve robustness to model variation in learned
control \cite{tobin2017domain,peng2018sim}.

This distinction determines how our evaluation should be read. Training
covers a declared family of load, harvesting, outage, and backhaul
conditions. We then freeze the policy and evaluate it on disjoint stochastic
traces drawn from that family. These experiments assess fixed-policy
generalization; they do not establish few-shot adaptation or robustness to
operating changes outside the specified family.

\subsection{Energy-Aware Flight and Safe Execution}
UAV service decisions share a finite battery with propulsion, communication,
and computation. Rotary-wing flight models therefore appear together with
offloading and trajectory decisions in UAV-MEC studies
\cite{liu2020uav,talukdar2026mappo}. Predictive UAV deployment also links
position to demand and mobility energy \cite{zhang2019predictive}. In this
work, immediate battery causality and a verified return/landing reserve are
hard execution constraints, while a virtual deficit queue represents the
declared long-run net-energy target. The safety controller checks complete
slot trajectories, rather than relying on a reward penalty to prevent
collision or battery exhaustion.

\subsection{Synthesis and Research Gap}
Table~\ref{tab:positioning} relates these research streams to the design
choices in this work. Although prior studies establish the individual value of
mobile edge servers, Lyapunov queue control, multi-agent learning, robust
training, and energy-aware flight, combining these ingredients is not a direct
composition. A routing action changes the arrival process of a downstream
queue; a processor or backhaul allocation changes both task delay and energy
consumption; and UAV motion changes the future channel, propulsion demand, and
return reserve. Treating any one of these effects outside the common state and
transition model can transfer congestion or energy cost to a stage that the
controller does not observe.

Three specific gaps motivate the proposed framework. First, much of the
offloading literature represents service through an instantaneous rate or an
aggregate delay, which does not reveal where unfinished work waits after it is
dispatched. An evaluation based on such a model may count assigned work as
successful even when upload, computation, synchronization, or backhaul service
remains incomplete. Second, standard MADDPG or MAPPO supplies a learning
mechanism but does not by itself enforce processor capacity, battery causality,
continuous-path separation, or a safe return reserve. Reward penalties alone
cannot guarantee these constraints at execution. Third, a policy that performs
well in a nominal simulator is not automatically reliable under load,
harvesting, outage, or backhaul changes. Meta-learning addresses rapid
adaptation, but it assumes that post-deployment data and gradient updates are
available; that assumption does not match the zero-update deployment studied
here.

Accordingly, the missing capability is a fixed multi-agent policy that observes
end-to-end congestion, coordinates the mixed decisions of heterogeneous
agents, respects hard physical constraints, and is trained and tested across a
declared family of operating regimes. The proposed framework addresses this
gap through persistent task-level queues, Lyapunov-derived physical and virtual
queue signals, hybrid CTDE control, and a deterministic feasibility layer. Its
evaluation freezes the trained actors and uses paired held-out traces so that
differences between controllers are not caused by different traffic or channel
realizations. This design supports a claim of robustness within the declared
domain family; it does not claim guaranteed performance for arbitrary
out-of-distribution conditions.

\begin{table*}[!t]
\caption{Positioning of the Proposed Framework Relative to the Main Research Streams}
\label{tab:positioning}
\centering
\footnotesize
\begin{tabular}{p{0.19\textwidth}p{0.34\textwidth}p{0.39\textwidth}}
\toprule
Research stream & Established capability & Gap addressed in this work\\
\midrule
UAV-assisted MEC~\cite{liu2020uav,xiong2021energy,guo2024multi} & Joint offloading, computing, radio allocation, and trajectory-aware service & Explicit vehicle--RSU--UAV--cloud stage queues, task records, branch synchronization, and finite backhaul\\
Queue-aware VEC~\cite{neely2010stochastic,mao2017dynamic,sun2026qtop} & Lyapunov-weighted offloading and UAV path planning & Transport and staging queues, weighted CPU sharing, censored unfinished tasks, and virtual net-energy accounting\\
MADDPG/MAPPO~\cite{lowe2017multi,schulman2017proximal,talukdar2026mappo} & Coordinated deterministic or stochastic multi-agent policies & Masked categorical routing combined with continuous split, CPU, and flight actions\\
Meta-RL~\cite{finn2017model,nichol2018reptile,yan2026rmaddpg} & Initialization for limited-update adaptation & Fixed zero-update deployment; no few-shot claim is assigned to the evaluated policy\\
Domain randomization~\cite{tobin2017domain,peng2018sim} & Robust control over a declared parameter distribution & Paired held-out evaluation across traffic, harvesting, outage, and backhaul regimes\\
Energy-aware UAV control~\cite{liu2020uav,talukdar2026mappo,zhang2019predictive} & Propulsion and computation energy models & Whole-system energy ledger, battery causality, collision avoidance, and verified return/landing reserve\\
\bottomrule
\end{tabular}
\end{table*}
